\subsection{Counting zeros of the return displacement}
\label{sec:direct-zero-counting}

We count distinct zeros of \(\Delta_p\) on an interval of central
heights \(h\). Multiplication by a nonvanishing factor or a regular
monotone change of height coordinate preserves their number. To relate
the critical points of the displacement to the endpoint slopes, put
\(m_\sigma=hI_{\sigma,h}>0\) and define
\begin{equation}
\label{eq:direct-completed-log-ratio}
 \mathcal L=\log\left[\Pi_p'(I_-)\frac{m_-}{m_+}\right],
 \qquad
 \Delta_p'(h)=I_+'(h)(\e^{\mathcal L}-1).
\end{equation}
We express the crossing directions in a logarithmic coordinate that
increases as the central height decreases:
\begin{equation}
\label{eq:direct-source-clock-convention}
 T=-c_T\log(h/h_{\rm ref}),\qquad
 \partial_T=-\frac{h}{c_T}\partial_h,\qquad
 m_\sigma=-c_T I_{\sigma,T},
\end{equation}
where \(c_T>0\) and \(h_{\rm ref}>0\) depend only on the
parameters and are held fixed when differentiating in \(h\). The
height-range covers are ordered by increasing \(T\), hence decreasing \(h\).
Changing these scales replaces \(T\) by \(c_1T+c_2\), with
\(c_1>0\) constant for fixed parameters. This preserves zero-crossing
directions and the sign of curvature at a critical point. We write
\(\mathcal L(T)\) for the same slope ratio in this coordinate.

By \eqref{eq:direct-completed-log-ratio}, a bound of \(n\) for the
zeros of \(\mathcal L\) on one interval of central heights gives at
most \(n+1\) zeros of \(\Delta_p\), by Rolle's theorem. This bound
is uniform in the additive constant of \(\Delta_p\). It remains to
combine estimates from adjacent height ranges without losing their
crossing directions. The next lemma also allows a zero at a strict
maximum.

\begin{lemma}[Ordered zero crossings]
\label{lem:direct-oriented-zero-count}
Let \(f\in C^2(J)\), where \(J\) is an interval. If every
critical point of \(f\) satisfies \(f''<0\), then \(f\) has
at most one critical point and at most two distinct zeros.
More generally, suppose \(J\) has an ordered cover by an initial,
middle, and final interval, with overlaps allowed and any of these
intervals possibly empty. Assume that
\begin{enumerate}[label=(\roman*)]
\item every zero in the initial interval has \(f'>0\);
\item every critical point in the middle interval has \(f''<0\);
\item every zero in the final interval has \(f'<0\).
\end{enumerate}
Then \(f\) has at most two distinct zeros on \(J\).
If every zero on \(J\) has the same strict derivative sign,
there is at most one. An interval on which \(f\) never vanishes
may be included in either end interval.
\end{lemma}

\begin{proof}
At a zero of \(f'\) with \(f''<0\), the derivative crosses
strictly from positive to negative. Two such crossings would require
an intervening zero of \(f'\) that is not a downward crossing.
Thus \(f\) is monotone or increases and then decreases. It has at
most two zeros; when there are two, an upward crossing is followed by
a downward crossing. A zero at the critical point is a strict maximum
and is the only zero. The same argument gives at most one zero when
\(f'\) has the same strict sign at every zero of \(f\).

For the ordered cover, choose cuts in the overlaps and assign each
shared endpoint to only one adjacent interval. In the initial interval
zeros cross upward; in the middle interval they occur in the order
just proved; in the final interval they cross downward. Two upward
crossings require an intervening downward crossing, and two downward
crossings require an intervening upward crossing. Either configuration
contradicts the permitted crossing directions and their order. A zero
at a middle maximum likewise excludes zeros in both end intervals.
Hence the only possible pair consists of an upward crossing followed
by a downward crossing.
\end{proof}

\begin{samepage}
\begin{corollary}[Ordered cover count]
\label{cor:ordered-cover-zero-count}
Let \(J\) be an interval of central heights on which \(\mathcal L\) is
\(C^2\), and let \(T\) be a logarithmic height coordinate as in
\eqref{eq:direct-source-clock-convention}. If \(J\) has an ordered cover,
in increasing \(T\), satisfying the hypotheses of
\cref{lem:direct-oriented-zero-count} for \(f=\mathcal L\), then
\(\Delta_p\) has at most three distinct zeros on \(J\). If every zero of
\(\mathcal L\) on \(J\) has the same strict crossing direction, then
\(\Delta_p\) has at most two.
\end{corollary}
\end{samepage}

\begin{proof}
By \cref{lem:direct-oriented-zero-count}, \(\mathcal L\) has at most two
zeros on \(J\), and at most one in the second case. Since
\(\Delta_T=-(m_+/c_T)(\e^{\mathcal L}-1)\) with \(m_+>0\), the critical
points of \(\Delta_p\) are the zeros of \(\mathcal L\). Rolle's theorem
gives the assertion.
\end{proof}

\begin{figure}[t]
\centering
\begin{tikzpicture}[x=1.05cm,y=.8cm,>=Stealth,font=\small]
  % logarithmic slope ratio
  \draw[->] (0,0)--(10.5,0) node[right] {$T$};
  \draw[->] (0,-1.25)--(0,1.75) node[above] {$\mathcal L$};
  \draw[thick] plot[smooth,tension=.7] coordinates
    {(0.2,-1) (1.5,-.6) (2.6,0) (3.7,.82) (5,1.18) (6.3,.82) (7.4,0)
     (8.5,-.6) (10,-.95)};
  \fill (2.6,0) circle (1.7pt) (7.4,0) circle (1.7pt) (5,1.18) circle (1.7pt);
  \node[above left] at (2.55,0) {upward};
  \node[above right] at (7.45,0) {downward};
  \node[above] at (5,1.2) {maximum};
  % ordered cover
  \draw[line width=2.4pt,gray!55] (0.2,-1.8)--(4.2,-1.8);
  \draw[line width=2.4pt,gray!55] (3.3,-2.3)--(6.9,-2.3);
  \draw[line width=2.4pt,gray!55] (6.1,-2.8)--(10,-2.8);
  \node[right] at (4.25,-1.8) {initial};
  \node[right] at (7.65,-2.3) {middle};
  \node[left] at (6.05,-2.8) {final};
  % displacement
  \draw[->] (0,-4.6)--(10.5,-4.6) node[right] {$T$};
  \draw[->] (0,-5.85)--(0,-3.4) node[above] {$\Delta_p$};
  \draw[thick,domain=0.3:9.8,samples=80] plot
    (\x,{-4.6+0.0814*((\x-5)*(\x-5)*(\x-5)/3-5.76*(\x-5))});
  \draw[dashed,gray] (2.6,0)--(2.6,-3.85);
  \draw[dashed,gray] (7.4,0)--(7.4,-5.35);
\end{tikzpicture}
\caption{Schematic of \cref{cor:ordered-cover-zero-count}, in increasing
\(T\), that is, decreasing central height: \(\mathcal L\) with an upward
zero, a strict maximum and a downward zero; the initial, middle and final
intervals of an ordered cover; and \(\Delta_p\), whose critical points
(dashed) are the zeros of \(\mathcal L\).}
\label{fig:ordered-cover}
\end{figure}


\Cref{fig:ordered-cover} shows the configuration allowed by
\cref{cor:ordered-cover-zero-count}. In the applications, thresholds
on the half-orbits define the pieces of the cover. The following
lemma uses the graph order \eqref{eq:central-graph-order-new} to
place these pieces in their order in \(T\).

\begin{lemma}[Order of central heights]
\label{lem:central-height-order}
Let \(h_1<h_2\) be central heights whose half-orbits of the same sign
\(\sigma\) are defined up to a common height \(z\). Then
\(|\eta_\sigma(z;h_1)|>|\eta_\sigma(z;h_2)|\), and at every common radius
the half-orbit through \(h_2\) has the larger height. Consequently, for
fixed \(z_e,R>0\), the central heights with \(h\ge z_e\) or
\(|\eta_+(z_e;h)|\le R\) form an interval that contains all larger
heights, that is, an initial interval in \(T\). For a fixed height
\(z_o\) and an increasing function \(g\), the central heights with
\(g(|\eta_+(z_o;h)|)\ge c_{\rm th}\) form a final interval in \(T\). The
intersection of such intervals with \(\cH_p\) is an interval.
\end{lemma}

\begin{proof}
The first assertion is \eqref{eq:central-graph-order-new}, since
\(|\eta_\sigma|=\sqrt{2U_\sigma}\). Each graph is strictly increasing in
\(z\): the first-zero argument \eqref{eq:graph-slope-first-zero}
gives \(zB-q>0\), using \(B+zB_z>0\) on the reduced normal-rate box.
Hence, at a common radius, the graph through the larger central
height has the larger height. The remaining assertions follow because
\(T=-c_T\log(h/h_{\rm ref})\) decreases with \(h\) and
\(\cH_p\) is an interval.
\end{proof}

\Cref{tab:ordered-covers} lists the ordered covers used in
\cref{sec:contact-positive-offsets,sec:recurring-ranges,sec:primary-regimes,sec:quadratic-regimes,sec:direct-negative-annulus-buffer,sec:direct-negative-root-regimes}.
The remaining regional statements reduce to
\cref{prop:weak-root-passage}.
\begin{table}[ht]
\centering
\small
\caption{Ordered covers of \(\cH_p\), in increasing \(T\). ``Up'' and ``down'' mean that every zero of \(\mathcal L\) in the
range crosses upward, respectively downward; ``max'' means that every
critical point of \(\mathcal L\) in the range is a strict maximum. The
bounds follow from \cref{cor:ordered-cover-zero-count} unless another
lemma is named.}
\label{tab:ordered-covers}
\begin{tabular}{@{}>{\raggedright\arraybackslash}p{\dimexpr.22\textwidth-8pt}
                >{\raggedright\arraybackslash}p{\dimexpr.66\textwidth-8pt}
                >{\raggedright\arraybackslash}p{\dimexpr.12\textwidth-8pt}@{}}
\hline
Proposition & Ranges and the property of \(\mathcal L\) & Zero bound for
  \(\Delta_p\)\\
\hline
\ref{prop:direct-bounded-contact}
  & large heights, \(\mathcal L<0\); compact heights, up, by
    \cref{lem:direct-contact-scattering-ratio,lem:direct-derivative-ratio-count};
    small heights, \(\mathcal L>0\)
  & 2\\
\ref{prop:direct-positive-offset}
  & heights above the saddle collar, \(\mathcal L<0\); the collar, up, by
    \cref{lem:one-logarithm-zero-count}
  & 2\\
\ref{prop:weak-root-passage}
  & early heights and the root interval, up
  & 2\\
\ref{prop:direct-primary-negative}
  & early heights, \(\mathcal L<0\); bounded radii, max, by
    \cref{lem:curvature-criterion}; large exits, \(\mathcal L<0\)
  & 3\\
\ref{prop:direct-positive-quadratic}, compact drift
  & early heights, up; the root interval, max, by
    \cref{lem:curvature-criterion}
  & 3\\
\ref{prop:direct-positive-quadratic}, small \(b\)
  & early heights and the root collar, up, by
    \cref{lem:one-logarithm-zero-count} on the collar
  & 2\\
\ref{prop:direct-buffered-negative}
  & early heights, up; annular exits, max, or \(\mathcal L<0\) for
    small \(b\); later exits, down
  & 3\\
\ref{prop:direct-fold-bounded}, \ref{prop:direct-negative-fold}
  & early heights, up; bounded outputs and escaped heights, max
  & 3\\
\ref{prop:direct-negative-separated}
  & early heights, up; the root interval, max
  & 3\\
\hline
\end{tabular}
\end{table}

Near a boundary root, taking the logarithm of a power-law derivative
reduces the slope comparison to a constant, one logarithm, and a
remainder. The next lemma uses the strict negative constant to control
the zeros even as the coefficient of the logarithm tends to zero.

\begin{samepage}
\begin{lemma}[One logarithm]
\label{lem:one-logarithm-zero-count}
For \(c_0,C_e,\vartheta>0\) there is \(x_c\in(0,1)\) with the following
property. Let \(G\in C^1((0,x_c))\) have the form
\[
 G(x)=c+\alpha\log x+e(x),\qquad c\le-c_0,\quad \alpha\in\R,\quad
 |e(x)|+|xe'(x)|\le C_e x^\vartheta .
\]
If \(\alpha\ge0\), then \(G<0\). In general \(xG'(x)<0\) at every zero of
\(G\), so \(G\) has at most one zero. If \(f\in C^1((0,x_c))\)
and \(f'\) has the sign of \(G\), then \(f\) has at most
two zeros.
\end{lemma}
\end{samepage}

\begin{proof}
Choose \(x_c\) so that \(C_e x^\vartheta<c_0/2\) and
\(C_e x^\vartheta|\log x|<c_0/4\) for \(0<x<x_c\). If \(\alpha\ge0\), then
\(G\le c+e<-c_0/2\). If \(\alpha<0\), every zero satisfies
\(|\alpha||\log x|=|c+e(x)|\ge c_0/2\), and hence
\(xG'(x)=\alpha+xe'(x)<-|\alpha|+|\alpha|/2<0\). All zeros are crossings
in the same direction, so there is at most one. Rolle's theorem gives the
bound for \(f\).
\end{proof}

\begin{remark}
\label{lem:direct-frozen-section-sign}
The auxiliary sections used in a local representation are held fixed
during differentiation. If each point of an interval \(J\) has a
neighborhood on which \(f\in C^1(J)\) is represented with auxiliary
sections fixed for that neighborhood, then a sign of \(f'\) proved in
these representations holds on all of \(J\).
Endpoint preparations formed from the same \(I_\pm\) and \(\Pi_p\) on
\(\widetilde W_\pm\) of \cref{lem:exterior-multiplier-new} represent the
same \(\mathcal L\) in \eqref{eq:direct-completed-log-ratio}, including
shared cuts.
\end{remark}

\subsection{Derivative-ratio comparisons on compact intervals}

For analytic functions, the classical Wronskian criterion
\cite[Lemma~2.3(b)]{GMV2011} identifies the extended complete
Chebyshev property by nonvanishing successive Wronskians.
For the three functions \(1,\phi,\psi\), the relevant identities are
\[
 W(1,\phi)=\phi',\qquad
 W(1,\phi,\psi)=(\phi')^2(\psi'/\phi')'.
\]
For the compact-interval comparisons below, we need this argument to
remain valid after adding a small \(C^2\) remainder. The next lemma
records the strict margins on the denominator and the derivative
of the ratio that make this possible.

\begin{lemma}[Stable derivative-ratio count]
\label{lem:direct-derivative-ratio-count}
On a fixed compact interval \(J\), let \(\phi,\psi\in C^2(J)\),
with \(\phi'\ge c_*>0\) and a uniform \(C^2\) bound on \(\phi\).
Let \(0<\lambda_0\le\lambda\le\lambda_1\). If
\[
 |(\psi'/\phi')'|\ge c_*
 \quad\hbox{on }J,
\]
then, for every constant \(D\) and every sufficiently small
\(C^2\) remainder \(R\), the function
\(g=D+\phi-\lambda\psi+R\) has at most two zeros on \(J\), and
\((g'/\phi')'\) has the sign of \(-(\psi'/\phi')'\) on \(J\).
If instead \(1-\lambda\psi'/\phi'\ge c_*\), then \(g'/\phi'>0\), and
\(g\) has at most one zero.
The remainder threshold depends only on the displayed uniform bounds.
\end{lemma}

\begin{proof}
Dividing the derivative by \(\phi'>0\) gives
\(g'/\phi'=1-\lambda\psi'/\phi'+R'/\phi'\).
The derivative of the remainder term is bounded by
\(C\|R\|_{C^2}\). In the first case, the continuous function
\((\psi'/\phi')'\) has one sign and
\(|\lambda(\psi'/\phi')'|\ge\lambda_0c_*\). Consequently,
\((g'/\phi')'=-\lambda(\psi'/\phi')'+O(\|R\|_{C^2})\) has the sign of
\(-(\psi'/\phi')'\) when the remainder is sufficiently small.
The quotient is strictly monotone, so Rolle's theorem gives at most
two zeros of \(g\). In the second case, the quotient remains positive,
which gives at most one zero.
\end{proof}
